% GENERATED FILE. Edit canonical lessons, then run npm run generate:books.
% canonical-source-hash: fnv1a64:6721cad1332b52a6
% canonical-lessons: GE-C206-wunde, GE-C206-pflaster, GE-C206-salbe, GE-C206-spritze, GE-C206-notfall

\chapter{A wound, A plaster, An ointment, An injection, An emergency}
\label{ch:ge-a-wound-a-plaster-an-ointment-an-injection-an-emergency}
\hlchaptermodality{206}

\begin{chapteropening}
\textbf{By the end of this chapter:} \emph{I can say a wound, a plaster, an ointment, an injection and an emergency.}

Complete the last lesson of chapter 206: I can say a wound, a plaster, an ointment, an injection and an emergency.
\end{chapteropening}

\section[die Wunde]{die Wunde — a wound}

\label{lesson:GE-C206-wunde}



\begin{quote}
Before the new one: say the German for an examination, then the German for an injury.
\end{quote}

\subsection*{You'll want to know: die Wunde}
\textbf{die Wunde} — \textquotedblleft{}a wound\textquotedblright{}.

\begin{cousinweb}
\textbf{Wunde} $\leftarrow$ Middle High German \emph{wunde} $\leftarrow$ Old High German \emph{wunta} — inherited, the same word as English \emph{wound}. (EtymDB 2.1)
\end{cousinweb}

\begin{grammarlens}[title={What you've built}]
One more word for this chapter.
\end{grammarlens}

\subsection*{Your turn}
\begin{itemize}
\raggedright
  \item \emph{Say it:} \emph{die Wunde}
  \item \emph{Say it:} \emph{die Wunde}, once more
  \item \emph{Say it:} say \emph{die Wunde}
  \item \emph{Recall:} say the German for an examination, then the German for an injury, then say \emph{die Wunde} again
\end{itemize}

\begin{tcolorbox}[breakable,colback=teal!4,colframe=teal!35!black,title={Before you move on}]
What does die Wunde mean? (\textquotedblleft{}A wound\textquotedblright{}.) Say it once more.
\end{tcolorbox}
\section[das Pflaster]{das Pflaster — a plaster}

\label{lesson:GE-C206-pflaster}



\begin{quote}
Before the new one: say the German for an injury, then the German for a wound.
\end{quote}

\subsection*{You'll want to know: das Pflaster}
\textbf{das Pflaster} — \textquotedblleft{}a plaster\textquotedblright{}.

\begin{cousinweb}
\textbf{Pflaster} $\leftarrow$ Old High German \emph{pflastar} $\leftarrow$ Latin \emph{emplastrum}, \textquotedblleft{}a salve spread on\textquotedblright{} (EtymDB 2.1) — borrowed, like English \emph{plaster}.
\end{cousinweb}

\begin{grammarlens}[title={What you've built}]
One more word for this chapter.
\end{grammarlens}

\subsection*{Your turn}
\begin{itemize}
\raggedright
  \item \emph{Say it:} \emph{das Pflaster}
  \item \emph{Say it:} \emph{das Pflaster}, once more
  \item \emph{Say it:} say \emph{das Pflaster}
  \item \emph{Recall:} say the German for an injury, then the German for a wound, then say \emph{das Pflaster} again
\end{itemize}

\begin{tcolorbox}[breakable,colback=teal!4,colframe=teal!35!black,title={Before you move on}]
What does das Pflaster mean? (\textquotedblleft{}A plaster\textquotedblright{}.) Say it once more.
\end{tcolorbox}
\section[die Salbe]{die Salbe — an ointment}

\label{lesson:GE-C206-salbe}



\begin{quote}
Before the new one: say the German for a wound, then the German for a plaster.
\end{quote}

\subsection*{You'll want to know: die Salbe}
\textbf{die Salbe} — \textquotedblleft{}an ointment\textquotedblright{}.

\begin{cousinweb}
\textbf{Salbe} $\leftarrow$ Middle High German \emph{salbe} $\leftarrow$ Old High German \emph{salba} — inherited, the same word as English \emph{salve}. (EtymDB 2.1)
\end{cousinweb}

\begin{grammarlens}[title={What you've built}]
One more word for this chapter.
\end{grammarlens}

\subsection*{Your turn}
\begin{itemize}
\raggedright
  \item \emph{Say it:} \emph{die Salbe}
  \item \emph{Say it:} \emph{die Salbe}, once more
  \item \emph{Say it:} say \emph{die Salbe}
  \item \emph{Recall:} say the German for a wound, then the German for a plaster, then say \emph{die Salbe} again
\end{itemize}

\begin{tcolorbox}[breakable,colback=teal!4,colframe=teal!35!black,title={Before you move on}]
What does die Salbe mean? (\textquotedblleft{}An ointment\textquotedblright{}.) Say it once more.
\end{tcolorbox}
\section[die Spritze]{die Spritze — an injection}

\label{lesson:GE-C206-spritze}



\begin{quote}
Before the new one: say the German for a plaster, then the German for an ointment.
\end{quote}

\subsection*{You'll want to know: die Spritze}
\textbf{die Spritze} — \textquotedblleft{}an injection\textquotedblright{}.

\begin{cousinweb}
\textbf{Spritze} comes from the verb \textquotedblleft{}to squirt, to spray\textquotedblright{} (EtymDB 2.1), shaped to sound like a spurt of liquid.
\end{cousinweb}

\begin{grammarlens}[title={What you've built}]
One more word for this chapter.
\end{grammarlens}

\subsection*{Your turn}
\begin{itemize}
\raggedright
  \item \emph{Say it:} \emph{die Spritze}
  \item \emph{Say it:} \emph{die Spritze}, once more
  \item \emph{Say it:} say \emph{die Spritze}
  \item \emph{Recall:} say the German for a plaster, then the German for an ointment, then say \emph{die Spritze} again
\end{itemize}

\begin{tcolorbox}[breakable,colback=teal!4,colframe=teal!35!black,title={Before you move on}]
What does die Spritze mean? (\textquotedblleft{}An injection\textquotedblright{}.) Say it once more.
\end{tcolorbox}
\section[der Notfall]{der Notfall — an emergency}

\label{lesson:GE-C206-notfall}



\begin{quote}
Before the new one: say the German for an ointment, then the German for an injection.
\end{quote}

\subsection*{You'll want to know: der Notfall}
\textbf{der Notfall} — \textquotedblleft{}an emergency\textquotedblright{}.

\begin{cousinweb}
\textbf{Notfall} is the word for \textquotedblleft{}need, distress\textquotedblright{}, Middle High German \emph{nōt}, a cousin of English \emph{need} (EtymDB 2.1), + the word for \textquotedblleft{}case, fall\textquotedblright{}: a case of distress.
\end{cousinweb}

\begin{grammarlens}[title={What you've built}]
One more word for this chapter.
\end{grammarlens}

\subsection*{Your turn}
\begin{itemize}
\raggedright
  \item \emph{Say it:} \emph{der Notfall}
  \item \emph{Say it:} \emph{der Notfall}, once more
  \item \emph{Say it:} say \emph{der Notfall}
  \item \emph{Recall:} say the German for an ointment, then the German for an injection, then say \emph{der Notfall} again
\end{itemize}

\begin{tcolorbox}[breakable,colback=teal!4,colframe=teal!35!black,title={Before you move on}]
What does der Notfall mean? (\textquotedblleft{}An emergency\textquotedblright{}.) Say it once more.
\end{tcolorbox}
